\chapter{先玩一个例子：为什么接近相等还不够}\label{sec:main}
\section{每一步到底允许做什么}
看这七个数：
\[(-3,0,0,0,1,1,1).
\]
每次选三个位置，把这三个数同时改成平均值。先选 $(-3,0,0)$，它们都变成 $-1$。现在有三个 $-1$、三个 $1$、一个 $0$。再选 $(-1,0,1)$，它们都变成零；这样重复三次，整排数就全为零。总共只做了四次平均。

这里的三个零是三个不同位置，不是一个数只写了一遍。以后我们用 $0^3$ 表示“三份零”，用 $a^m$ 表示“$m$ 个位置都取值 $a$”。在这种配置记号里，上标表示份数。

现在把“三个”换成一个素数 $p$，把总位置数记为 $n$。两个字母管不同的事：$p$ 是一次操作选多少个，$n$ 是总共能使用多少个位置。所有操作都在原来的 $n$ 个位置上进行。

\section{终点其实已经确定了}
一次操作把 $p$ 个数换成 $p$ 份平均值，总和没有变化。因此，假如最终全相等，公共值必然是
\[A=\frac{x_1+\cdots+x_n}{n}.
\]
我们不必猜终点，只需判断能不能到达这个确定的 $A$。

“越来越接近 $A$”和“某一步恰好全等于 $A$”之间有很大差别。数列 $1,1/2,1/4,\ldots$ 越来越小，却没有一项是零。平均确实会让差异变小，但仅凭这一点还不能回答有限步的问题。

\section{先把题目化成整数归零}
把所有坐标减去 $A$，终点就变成零。再同时乘一个数清除分母，最后除去坐标共同的最大公因数。这不改变成功的位置选择序列，因为取平均与整体平移、整体缩放交换。

例如输入 $(1,2,4)$ 的均值为 $7/3$。减去均值得 $(-4/3,-1/3,5/3)$，再乘 $3$ 得
\[X=(-4,-1,5).
\]
三个坐标的最大公因数是 $1$，不用再约分。我们把这样的整数代表叫作\textbf{本原中心化代表}：“中心化”表示总和为零，“本原”表示已经约掉共同整数因子。

一般写成
\[X=c(x-A\mathbf1)\in\Z^n,\qquad \sum_iX_i=0,\qquad\gcd_iX_i=1.
\]
对非恒定输入，这个代表只差一个整体正负号。现在定义
\begin{equation}\label{eq:G}
G(X)=\gcd_{i<j}|X_i-X_j|.
\end{equation}
这是\emph{差的}最大公因数，与坐标自己的最大公因数不同。上例坐标本原，但三个差为 $3,9,6$，所以 $G=3$。下一讲会解释，为什么恰好该看这些差。

\section{我们要证明的答案}
\InsertProof{thm:prime}
这里 $p^0=1$ 也允许。这个定理不是说位置足够多以后所有配置都能平均；它说此时障碍恰好由 $G$ 是否含非 $p$ 因子决定。
\InsertProof{01_problem-unit1}
\InsertProof{cor:threshold}
“最终阈值”是从某个位置数开始，\emph{每一个更大的位置数}都满足判据。例如 $n=p$ 时全体平均一次就完成，但 $n=2p$ 反而存在 $G=1$ 的失败例子。因此不能把“第一个成功的维数”和“以后永远成立的界”混为一谈。

\section{后面的证明怎样走}
先证明有非 $p$ 障碍必然失败。这部分很短。真正长的是反方向：没有这些障碍，为什么一定能选出一条成功路径？

位置较多时，我们让平方和作为非负整数不断下降；这需要精心挑选平均组。位置紧张时，我们把整排数整理成三个等值块，用两个参数记录；再用几次平均组成一个矩阵动作，把参数送到容易收尾的位置。二者最后都会遇到一个共同的好消息：一旦选出和为零的 $p$ 个数，接下来已有独立的收尾方法。

你现在只需记住这三件事：目标可以化成零；$G$ 负责判断最初的算术障碍；充分性要同时解决“下一步怎么选”和“为什么有限步结束”。接下来逐件说明。
